% ============================================================
% §6 模型评价与推广 (Stage 07 产物) — 精简版
% 竞赛: 电工杯 B题 (2026)
% 详细版见附录E
% ============================================================

\section{模型评价与推广}

\subsection{模型优点}

\begin{enumerate}[label=\textbf{优点\arabic*：}, leftmargin=*, itemsep=4pt]

    \item \textbf{双模型三角校验确立预测可信度边界}。增生型Markov链与Leslie矩阵两种独立范式对第5年末60+人口的预测偏差仅1.7\%，Leslie主特征值$\lambda_1=1.0198$与Markov隐含年均增长率2.0\%精确吻合，为需求预测提供了数理交叉验证的"黄金标准"。

    \item \textbf{MILP-MCLP-ES精确线性化框架保留全局最优性}。针对$S_2(u_i)\cdot x_{ij}$非线性乘积，以McCormick包络——4条线性约束——等价替换；$S_1$与$S_2$的分段常数函数——分别为4段与5段——以Big-M方法精确编码。将原始混合整数非线性规划等价转化为MILP，CBC分支定界可在120秒内收敛至gap$<1\%$的确定解。

    \item \textbf{词典序目标优雅解耦S3满意度与营收的多目标冲突}。优先最大化加权$S_3$，仅以$\varepsilon=10^{-4}$弱激励保留营收方向，确保全$S_3=1.000$平价区间的同时识别出唯一的紧约束变量——上门护理降价57.9\%。该结构为"公益优先、微利运营"提供了可复现的数学编码范式。

    \item \textbf{模型输出具备直接工程可操作性}。Q2选址方案可转化为招标标段划分，Q3定价已核算至0.01元精度，Q4预算相变分析为"追加多少投资可实现满意度跃升"提供了精确的边际锚点——13万元触发$S_2:0.600\to1.000$。

\end{enumerate}

\subsection{模型局限与改进方向}

\begin{enumerate}[label=\textbf{局限\arabic*：}, leftmargin=*, itemsep=4pt]

    \item \textbf{序贯求解非联合最优}。Q1$\to$Q2$\to$Q3的串行依赖链虽保证各阶段独立最优，但非全局联合最优。Benders分解或列-约束生成可实现预测-选址-定价联合优化，预期目标值提升3--5\%，但开发时间+15--20h对竞赛窗口不可行，更适合赛后学术扩展。

    \item \textbf{Markov转移概率的平稳性假设}。年际转移概率恒定忽视了年龄队列效应——高龄老人退化加速——与政策干预效应。时变转移矩阵$\mathbf{A}(t)$或贝叶斯层次模型可捕捉时变性，但需3--5年纵向追踪数据校准。

    \item \textbf{$S_2$全员触底的满意度悬崖}。全覆盖方案下三站利用率均超92\%，$S_2$全员锁定0.600。引入软容量约束——允许5--10\%超载但$S_2$线性惩罚——或85\%容量缓冲系数，可释放$S_2:0.600\to0.720$的跃迁，但需追加约13万元容量补偿投资。

    \item \textbf{确定性优化对参数敏感的脆弱性}。Q2 MILP对随机种子（分支策略）敏感，多起点随机重启——10次独立求解取最优——可将期望目标值提升3--8\%，代价为计算时间$\times 10$。分布鲁棒优化以Wasserstein模糊集替代确定性需求，可从根本上消除种子敏感性。

\end{enumerate}

\subsection{模型推广}

本框架的"MCLP覆盖+分段满意度+三层交叉补贴"核心方法论可迁移至以下公共设施优化场景：

\begin{enumerate}[label=\textbf{(\arabic*)}, leftmargin=*, itemsep=3pt]
    \item \textbf{分级诊疗站点选址}。"15分钟医疗服务圈"与MCLP覆盖目标同构，老年人口类型$\to$疾病严重度分级，服务项目$\to$诊疗科目，直接对接国家卫健委2027年75\%基层签约覆盖率政策目标。
    \item \textbf{冷链前置仓网络规划}。三类老人$\to$三温区商品品类——冷冻、冷藏与常温，配送时效满意度替代距离满意度，选址-容量-定价联合优化可缓解前置仓亏损率30--40\%的行业痛点。
    \item \textbf{新能源充电桩布局}。消除"充电荒漠"与MCLP全覆盖目标精确对齐，补贴帽机制可防止"站点越建越大、补贴效率越来越低"的财政资源错配。
\end{enumerate}
